% 张三慧V2第18章PDF物理页156源思考题18.1--18.11；原图18.18/18.19需复核。
\section*{第18章\quad 磁力：思考题}
\paragraph{思考题 18.1}
若已知电荷q以速率v运动经过某点时在不同速度方向受到的最大磁力，怎样定义该点的磁感应强度B？

\paragraph{思考题 18.2}
高能带电宇宙射线进入地球附近磁场时，为什么靠近两极的进入概率可高于赤道地区？

\paragraph{思考题 18.3}
若欲在赤道附近向西或向东发射质子，使其尽量沿赤道附近保持运动，地磁场力应如何判断？

\paragraph{思考题 18.4}
赤道附近地磁场近似水平向北，电场竖直向地面。电子要同时避免电场与磁场造成的横向偏转，应沿哪一水平东西方向发射？

\paragraph{思考题 18.5}
能否利用磁场对带电粒子做功而提高其动能？

\paragraph{思考题 18.6}
带电粒子在磁场由弱区进入强区作螺旋运动，平行和垂直磁场的速度分量、动能与磁矩会怎样变化？说明绝热条件。

\paragraph{思考题 18.7}
长直载流导线周围是不均匀的环形磁场，正电粒子沿着导线大致平行射入，受力方向和轨迹为何不一定是简单圆周？

\paragraph{思考题 18.8}
相互垂直的均匀电场E与磁场B构成速度选择器，怎样选择粒子才能无横向偏转通过？粒子电荷正负是否影响选中速率？

\paragraph{思考题 18.9}
磁场方向与电流方向给定时，载流导体受到的安培力方向是否由载流电子或正电荷的微观载流种类决定？霍尔电压极性又如何？

\paragraph{思考题 18.10}
原图18.18电子和正电子在垂直纸面的同一磁场中形成两条相反旋向的螺旋云室轨迹，怎样由曲率方向和半径变化辨别各电荷？

\paragraph{思考题 18.11}
图18.19中均匀电场向上、均匀磁场垂直纸面向外，一颗正质子在竖直平面内沿圆弧和尖点形轨迹运动。其速率如何随所在位置和x方向往返变化？
